\documentclass[aps,prb,preprint,superscriptaddress,floatfix,nofootinbib]{revtex4-2}

\usepackage{amsmath,amssymb,booktabs,array}
\usepackage[colorlinks=true,citecolor=blue,linkcolor=blue,urlcolor=blue]{hyperref}

\newcommand{\code}[1]{\texttt{#1}}

\begin{document}

\title{Supplemental Material for ``Interband Pairing Signatures in the Superfluid Response of Magic-Angle Twisted Bilayer Graphene''}

\author{Jian Zhou}
\affiliation{JZ Institute of Science, Hong Kong, China}
\email{jack@jzis.org}

\date{June 9, 2026}

\maketitle

\section{Scope of the Supplemental Material}

This Supplemental Material documents the executable numerical workflow behind
the PRB manuscript.  It is organized as a reproducibility and self-audit record:
the Bistritzer-MacDonald (BM) model conventions~\cite{bistritzer_2011}, the
orbital-projected pairing construction, the finite-band response formula, the
figure and table provenance, and the historical PRB benchmark reconstruction
are recorded separately.

The main text makes a conservative mechanism claim: an orbital-projected
interband-pairing direction produces a reproducible, normalization-conditioned
change in the normalized finite-band superfluid response.  The present
Supplemental Material does not promote the raw response units to final
experimental stiffness values.  This choice is motivated by the distinction
between robust quantum-geometric mechanism diagnostics~\cite{peotta_torma_2015,hu_2019,julku_2020,xie_2020_prl}
and direct quantitative comparison with recent stiffness
experiments~\cite{tian_2023,tanaka_2025}.

\section{BM Model and Numerical Parameters}

The normal-state calculation uses the single-valley BM continuum Hamiltonian
with orbital basis
\begin{equation}
  (tA,tB,bA,bB),
\end{equation}
where $t,b$ label layers and $A,B$ label sublattices.  Unless explicitly stated
otherwise, the production mechanism scans use the parameters in
Table~\ref{tab:sm_parameters}.

\begin{table}[t]
\caption{Default parameters used for the mechanism scans.}
\label{tab:sm_parameters}
\begin{ruledtabular}
\begin{tabular}{ll}
quantity & value\\
\hline
twist angle & $\theta=1.05^\circ$\\
Dirac velocity & $v_F=2.135~\mathrm{eV}\,\text{\AA}$\\
AA tunneling & $w_0=87.2~\mathrm{meV}$\\
AB tunneling & $w_1=109.0~\mathrm{meV}$\\
reciprocal cutoff & $N_{\rm shell}=3$\\
retained bands & $n_{\rm keep}=6$\\
momentum grid & $nk=7$ for dense maps, $nk=9$ for key-point validation\\
gap scale & $\Delta_0=1~\mathrm{meV}$\\
velocity finite difference & $dk=10^{-6}~\text{\AA}^{-1}$\\
\end{tabular}
\end{ruledtabular}
\end{table}

The moire unit-cell area and reciprocal scale used in the unit audit are
\begin{align}
  A_M &= 1.5606\times 10^4~\text{\AA}^{2}, &
  G_M &= 0.054047~\text{\AA}^{-1}.
\end{align}
The raw response conversion used in the audit is
\begin{equation}
  D[\mathrm{eV}\,\text{\AA}^{2}]
  =
  \frac{D_{\rm raw}[\mathrm{meV}\,\text{\AA}^{2}]}{1000},
\end{equation}
with all remaining spin, valley, or convention factors kept explicit rather
than absorbed into a hidden prefactor.

\section{Orbital-Projected Pairing Construction}

The pairing direction is specified first in the layer-sublattice orbital basis,
\begin{equation}
  \Delta_{\rm orb}(\eta)
  =
  \Delta_0\,\mathcal{N}_{\eta}(M_0+\eta M_1),
\end{equation}
with
\begin{equation}
  M_0=\tau_0\sigma_0,\qquad
  M_1=\tau_x\sigma_x.
\end{equation}
The two normalization controls used in the manuscript are
\begin{align}
  &\text{fixed Frobenius:}
  &&\|M_0+\eta M_1\|_F=\|M_0\|_F,\\
  &\text{fixed }\Delta_0:
  &&\mathcal{N}_{\eta}=1.
\end{align}

The band-basis pairing matrix is
\begin{equation}
  \Delta_{\rm band}(\mathbf{k})
  =
  U^\dagger(\mathbf{k})\Delta_{\rm orb}U^*(-\mathbf{k}).
\end{equation}
The working valley convention in the present diagnostic is the time-reversal
sewn proxy
\begin{equation}
  U_{-}(-\mathbf{k}) = U_{+}^*(\mathbf{k}).
\end{equation}
This convention is useful for isolating an orbital-defined interband-pairing
signature, but the readiness audit identifies a fully documented two-valley
implementation as a submission-critical improvement.

The projected interband-pairing weight is
\begin{equation}
  W_{\rm inter}
  =
  \frac{\sum_{n\ne m}|\Delta_{nm}|^2}
       {\sum_{n,m}|\Delta_{nm}|^2}.
\end{equation}
For the dense $nk=7$, $n_{\rm keep}=6$ scan, $W_{\rm inter}\simeq 0.108$ at
$\eta=1$.

\section{Finite-Band Response and Channel Split}

For a retained band subspace, the finite-band BdG Hamiltonian is
\begin{equation}
  H_{\rm BdG}(\mathbf{k}) =
  \begin{pmatrix}
  \varepsilon(\mathbf{k})-\mu & \Delta_{\rm band}(\mathbf{k})\\
  \Delta^\dagger_{\rm band}(\mathbf{k}) & -[\varepsilon(\mathbf{k})-\mu]
  \end{pmatrix}.
\end{equation}
The zero-temperature static current-current response is evaluated as
\begin{equation}
  \Lambda_{\alpha\beta}
  =
  \sum_{a\ne b}
  \frac{
  \langle a|J_\alpha|b\rangle
  \langle b|J_\beta|a\rangle
  (f_a-f_b)}
  {E_b-E_a}.
\end{equation}

The band-basis velocity is split into diagonal and off-diagonal parts,
\begin{equation}
  v_\alpha^{\rm band}
  =
  v_\alpha^{\rm intra}+v_\alpha^{\rm inter}.
\end{equation}
This produces conventional, geometric, and cross-channel response diagnostics.
The primary normalized response map uses
\begin{equation}
  \delta D_{\rm iso}(\mu,\eta)
  =
  \frac{D_{\rm iso}(\mu,\eta)}
       {D_{\rm iso}(\mu,0)} - 1,
  \qquad
  D_{\rm iso}=\frac{D_{xx}+D_{yy}}{2}.
\end{equation}

\section{Figure and Table Provenance}

Table~\ref{tab:sm_provenance} summarizes the current figure and table
provenance.  The full machine-readable record is maintained in
\code{notes/figure\_table\_provenance\_2026-06-09.md}.

\begin{table*}[t]
\caption{Manuscript artifact provenance.  CSV values used in the manuscript
tables are checked by \code{scripts/verify\_manuscript\_tables.py}.}
\label{tab:sm_provenance}
\begin{ruledtabular}
\begin{tabular}{llll}
artifact & data source & script & check\\
\hline
Frobenius heatmap & dense scan CSV & heatmap plotter & figure file exists\\
fixed-$\Delta_0$ heatmap & dense scan CSV & heatmap plotter & figure file exists\\
dense response table & eta-summary CSV & eta summarizer & table verifier\\
$nk=9$ validation table & key-point summary CSV & eta summarizer & table verifier\\
PRB audit table & reconstruction CSVs & audit scripts & table verifier\\
\end{tabular}
\end{ruledtabular}
\end{table*}

The dense response data are reproduced by
\begin{verbatim}
python3 scripts/run_mu_response_scan.py \
  --nk 7 --n-shell 3 --n-keep 6 \
  --mus -5 -4 -3 -2 -1 0 1 2 3 4 5 \
  --etas 0 0.25 0.5 0.75 1 \
  --output data/processed/mu_eta_response_scan_nk7_nkeep6.csv
\end{verbatim}
and summarized by
\begin{verbatim}
python3 scripts/summarize_eta_effect.py \
  data/processed/mu_eta_response_scan_nk7_nkeep6.csv \
  --output data/processed/mu_eta_response_scan_nk7_nkeep6_eta1_summary.csv
\end{verbatim}
The table verifier is
\begin{verbatim}
python3 scripts/verify_manuscript_tables.py
\end{verbatim}

\section{Key-Point Momentum-Grid Convergence}

The dense maps in the main text use $nk=7$, while selected key points were
checked at $nk=9$, $nk=11$, and $nk=13$.  The higher-grid scans use
\begin{verbatim}
python3 scripts/run_mu_response_scan.py \
  --nk 11 --n-shell 3 --n-keep 6 \
  --mus -4 0 2 4 --etas 0 0.5 1 \
  --output data/processed/mu_response_scan_nk11_nkeep6_keypoints.csv

python3 scripts/run_mu_response_scan.py \
  --nk 13 --n-shell 3 --n-keep 6 \
  --mus -4 0 2 4 --etas 0 0.5 1 \
  --output data/processed/mu_response_scan_nk13_nkeep6_keypoints.csv
\end{verbatim}
and is summarized with \code{scripts/summarize\_eta\_effect.py}.

Table~\ref{tab:sm_nk_convergence} compares the normalized $\eta=1$ response
relative to $\eta=0$.  The fixed-Frobenius response remains positive at all
four key chemical potentials.  The fixed-$\Delta_0$ response remains small and
sign-sensitive, consistent with the main-text interpretation that the
interband-pairing signal is normalization-conditioned.

\begin{table}[t]
\caption{Key-point momentum-grid check for
$D_{\rm iso}(\eta=1)/D_{\rm iso}(\eta=0)-1$.}
\label{tab:sm_nk_convergence}
\begin{ruledtabular}
\begin{tabular}{lrrrr}
mode, $\mu$ (meV) & $nk=9$ & $nk=11$ & $nk=13$ & $13-9$\\
\hline
fixed $\Delta_0$, -4 & -0.0068 & -0.0061 & +0.0122 & +0.0191\\
fixed $\Delta_0$,  0 & +0.0103 & +0.0080 & +0.0196 & +0.0093\\
fixed $\Delta_0$,  2 & -0.0029 & +0.0118 & +0.0024 & +0.0053\\
fixed $\Delta_0$,  4 & -0.0017 & -0.0003 & +0.0011 & +0.0029\\
fixed Frobenius, -4 & +0.0604 & +0.0617 & +0.0872 & +0.0267\\
fixed Frobenius,  0 & +0.0863 & +0.0796 & +0.1013 & +0.0150\\
fixed Frobenius,  2 & +0.0643 & +0.0927 & +0.0805 & +0.0162\\
fixed Frobenius,  4 & +0.0605 & +0.0608 & +0.0772 & +0.0166\\
\end{tabular}
\end{ruledtabular}
\end{table}

Across the three key-point grids the fixed-Frobenius response remains positive,
with the largest $nk=13-nk=9$ shift about 2.67 percentage points.  Therefore
the current manuscript claims sign and order-of-magnitude stability of the
normalized mechanism signal, not a strict continuum-limit extrapolation of the
tabulated values.

\section{Historical PRB Benchmark Reconstruction}

The production interband-pairing scan and the historical PRB table
reconstruction are kept separate.  The best simple unified reconstruction route
for the old uniform-$s$ benchmark is
\begin{equation}
  D_{\rm conv}=2D_{\rm intra,\tau_z},
  \qquad
  D_{\rm geom}=D_{\rm inter,\tau_z}.
\end{equation}
For $n_{\rm keep}=6$, this gives $D_{\rm total}=126.66~\mathrm{eV}\,\text{\AA}^2$
versus the historical target $129.30~\mathrm{eV}\,\text{\AA}^2$, and
$D_{\rm geom}=75.22~\mathrm{eV}\,\text{\AA}^2$ versus
$75.30~\mathrm{eV}\,\text{\AA}^2$.

The old $n_{\rm keep}=2$ endpoint is more convention-sensitive.  The closest
audited reconstruction combines a Gamma-centered mesh, central-pair band
selection, an all-$\tau_z$ current convention, a factor of two in the intraband
term, and a full curvature-like conventional correction.  This gives
\begin{equation}
  D_{\rm total}=66.18~\mathrm{eV}\,\text{\AA}^2
  \quad
  \text{versus}
  \quad
  67.50~\mathrm{eV}\,\text{\AA}^2.
\end{equation}
Because this full-curvature correction does not work uniformly at
$n_{\rm keep}=6$, the reconstruction is treated as a historical audit rather
than a production convention.

\section{Current Limitations and Next Validation Gates}

The current submission-critical gaps are:
\begin{enumerate}
  \item a systematic continuum-grid extrapolation beyond the present selected
  $nk=9$, $nk=11$, and $nk=13$ key-point checks;
  \item a fully documented two-valley sewing convention or an explicitly
  limited diagnostic-model framing;
  \item a final absolute stiffness convention before any direct experimental
  comparison;
  \item a calibrated filling variable replacing the retained-band proxy
  $\nu_{\rm proxy}$.
\end{enumerate}

These limitations do not invalidate the normalized-response mechanism claim,
but they prevent the current draft from being treated as a final PRB submission
without further convergence and convention work.

\begin{acknowledgments}
The author acknowledges use of open-source scientific computing tools
including NumPy and Matplotlib.
\end{acknowledgments}

\section*{Declarations}

\textbf{Funding}: Not applicable.

\textbf{Conflict of interest}: The author declares no conflict of interest.

\textbf{Data availability}: All data and code supporting this work are available
from the corresponding author upon reasonable request.

\textbf{Ethics approval}: Not applicable.

\bibliography{references}

\end{document}
